\chapter*{実装}

TOTPの計算方法も、Luaをnginxで扱う方法もわかったので、実装しましょう。
ユーザの登録機能は仕様範囲外とします。

\section*{ユーザーデータベース}
ユーザー名と鍵がPostgreSQLのusersテーブルに格納されているとします。
鍵長は10bitで、鍵は16進20桁のASCIIで格納しましょう。

\begin{lstlisting}[caption=SQL Schema,label=list:schema]
CREATE TABLE users (
	name text not null,
	secret varchar(20)
);
CREATE INDEX TO idxname ON users (name);
\end{lstlisting}

データベーススキーマをListing \ref{list:schema}に示します。
ここでは、Key-Value Storeっぽく使ってみたかったので、あえて正規化を
していません。まぁする必要もないですが。

Postgresからsecretを取得するために、ドライバをロードしておきましょう。
nginx.confにinit\_by\_luaを追記記述しましょう。
init\_by\_lua は、nginxの設定ファイルをロードする時に実行され、
Lua インタープリタの初期化に使います。
PostgreSQLのアクセスには、LuaSQLを使用しました。

\begin{lstlisting}[caption=Luaインタープリタの初期化,label=list:init]
init_by_lua '
        driver = require "luasql.postgres"
        env = assert(driver.postgres())
        ffi = require("ffi")
        ffi.cdef[[void get_token(const char *key, char *result, uint32_t len)]]
        otp = ffi.load("/usr/local/lib/libtotp.so")
';
server {
           ....
\end{lstlisting}

また、4行目以降はListing \ref{list:gettoken} のコードを含む
トークンを計算するためのライブラリをロードしています。

Listing \ref{list:getsecret}は、PostgreSQLからsecretを取得する関数です。
SQL Injectionを防ぐため、Prepared statementを使用しています。

\begin{lstlisting}[caption=データベースからsecretを取得する関数,label=list:getsecret]
function get_secret(name)
  local con=assert(env:connect("dbname=otp user=user"))
  con:execute("PREPARE secret (text) as SELECT secret from users where name=$1;")
  local cur=con:execute(string.format([[EXECUTE secret('%s');]],name))
  local row = cur:fetch({},"a")
  local secret
  if not row then
    secret = nil
  else
    secret = row.secret
  end
  cur:close()
  con:close()

  return secret
end
\end{lstlisting}

\section*{Cライブラリを用いたトークンの計算}
Listing \ref{list:token}は、FFIを使ってC++で書いたライブラリを
呼び出すところです。2行目で、結果を格納するchar型のバッファを確保して
Listing \ref{list:init}の6行目でロードしたライブラリのget\_token関数に
渡しています。
そして、char[64]型をLua stringに変換するためにffi.string関数を使用しました。

\begin{lstlisting}[caption=Tokenを計算,label=list:token]
function get_token(secret)
  local token = ffi.new("char[64]")
  otp.get_token(secret, token, 64)
  return ffi.string(token)
end
\end{lstlisting}

\section*{認証}
Listing \ref{list:auth}には、認証部分のコードを示します。
POSTメソッドで、ユーザ名とトークンが入力されると
DBにアクセスしてsecretを取得し、tokenを計算しています。
計算したtokenが入力されたtokenと一致したら、ngx.exit関数を
用いて、ステータスコード200を、ユーザーデータベースにユーザー名が
登録されていなかったり、トークンが一致しなかったら、401 Unauthorizedを
返却します。

\begin{lstlisting}[caption=SQL Schema,label=list:auth]
ngx.req.read_body()
local args, err = ngx.req.get_post_args()
if not args then
  ngx.exit(ngx.HTTP_UNAUTHORIZED)
end

local user = args.user
local secret = get_secret(user)
if not secret then
  ngx.exit(ngx.HTTP_UNAUTHORIZED)
end

if(args.token ~= get_token(secret)) then
  ngx.exit(ngx.HTTP_UNAUTHORIZED)
end

ngx.exit(ngx.HTTP_OK)
\end{lstlisting}

\section*{QRコードの作成}
Google Authenticatorは、QRコードでsecretを設定するのが一般的です。
secretを生成し、QRコードを生成してみましょう。
まず乱数源から10bit取り出して、BASE32エンコードします。
Listing \ref{list:secret}は、/dev/urandomから10bitの乱数を取り出して
BASE32エンコードしています。5行目の文字列が生成されたsecretです。

\begin{lstlisting}[caption=secret生成の例,label=list:secret]
$ dd if=/dev/urandom bs=10 count=1|base32
1+0 レコード入力
1+0 レコード出力
10 bytes copied, 0.000120597 s, 82.9 kB/s
NDS4MMILSEND5DZZ
\end{lstlisting}

次にQRコードを生成します。
Linuxを使っている方は、qrencodeコマンドを使用するのが楽です。
\begin{lstlisting}[caption=QRコード生成の例,label=list:qr]
$ qrencode -o qrcode.png "otpauth://totp/Example:alice@example.com?secret=NDS4MMILSEND5DZZ&issuer=Example"
\end{lstlisting}
とコマンドラインで実行すると、qrcode.pngというファイル名で、図\ref{fig:qr}
のようなQRコードが生成されます。

\begin{figure}
  \begin{center}
    \includegraphics{qrcode.eps} \label{fig:qr}
  \end{center}
  \caption{生成されたQRコード}
\end{figure}
%
Listing \ref{list:qr}でqrencodeに渡した文字列の、secretがトークンの生成に
使用する、BASE32エンコードされたsecretです。

次に生成したsecretをデータベースに登録します。
\begin{lstlisting}[caption=データベースへの登録,label=list:insert]
$ echo NDS4MMILSEND5DZZ | base32 -d |hexdump
0000000 e568 31c6 910b 3e1a 398f               
000000a
$ psql -c "INSERT INTO users (name,secret) VALUES('alice@example.com','68e5c6310b911a3e8f39');" otp
\end{lstlisting}
Listing \ref{list:insert}の1行目で生成したsecretを16進数に変換しています。
4行目は、その変換した値をpsqlコマンドでデータベースにINSERTしているのですが
注意して欲しいのは、登録するsecretがビッグエンディアンになっています。
もちろん、CPUがビッグエンディアンであるコンピュータを使用されていれば、
2行目に表示された通り入力します。

\section*{Deploy}
Luaのコードを配置しましょう。作成したLuaのプログラムをnginx.confに
埋め込んでも良いのですが、行数も多くなりましたので、JITコンパイルした
ファイルを指定することにしましょう。これにより、設定ファイルの可読性が
上がるだけでなく、実効速度も向上します。

JITコンパイルは以下ように行います。
\begin{lstlisting}[caption=JIT Compile,label=list:compile]
$ luajit -bg totp.lua totp.luac
\end{lstlisting}
%
totp.luaをコンパイルして、totp.luacを作成します。これを設定ファイルで
指定します。locationブロックのaccess\_by\_lua\_fileに指定します。

\begin{lstlisting}[caption=コンパイルしたファイルを設定,label=list:config]
server{
          ......
  location /auth{
     default_type text/plain;
     content_by_lua_file '/usr/local/lib/totp/totp.luac';
  }
          .....
\end{lstlisting}


\section*{テスト}
テストは、curlコマンドでPOSTリクエストを送るのが良いでしょう
\begin{lstlisting}[caption=テスト,label=list:test]
$ curl -v -X POST -d "user=alice" -d"token:123456" http://localhost/auth 
\end{lstlisting}
これで、ステータスコード200が返ってくれば認証成功、401が返ってくれば
認証失敗です。

